% introduction.tex — Paper 6, Section I: Introduction
% ASSERT_CONVENTION: natural_units=natural, metric_signature=mostly_plus, entropy_base=nats, modular_hamiltonian=K_minus_ln_rho, coupling_convention=H_sum_hxy, state_normalization=Tr_rho_1
\section{Introduction}
\label{sec:introduction}

Self-modeling---the condition that a physical system contains a faithful
model of itself, updated through its boundary---forces the algebraic
structure of quantum mechanics~\cite{Paper5}. We show that it also forces
general relativity. Specifically, Einstein's field equations emerge as the
leading-order infrared effective description of a lattice of self-modeling
systems in the Wilsonian continuum limit.

The idea that spacetime geometry is encoded in quantum entanglement has a
rich history. Van Raamsdonk argued that entanglement is the glue that
holds spacetime together~\cite{VanRaamsdonk2010}. Jacobson showed that the
Einstein equation can be derived as an equation of state from
thermodynamic considerations at local causal
horizons~\cite{Jacobson1995}, and later sharpened this to an entanglement
equilibrium argument requiring only area-law entanglement entropy, the
entanglement first law, and the maximal vacuum entanglement hypothesis
(MVEH)~\cite{Jacobson2016}. Lashkari, McDermott, and Van Raamsdonk
demonstrated that the entanglement first law in holographic conformal
field theories implies the linearized Einstein equation~\cite{LMVR2014},
and Faulkner~\textit{et~al.}\ extended this to the full nonlinear
equations in the holographic setting~\cite{Faulkner2014}. Cao, Carroll,
and Michalakis showed that spatial geometry can be recovered from the
entanglement structure of a finite-dimensional Hilbert
space~\cite{CCM2017}.

These results leave open a central question: \emph{what ultraviolet
completion produces the entanglement structure that Jacobson's argument
requires?} The holographic results of~\cite{LMVR2014,Faulkner2014} assume
AdS/CFT; the Jacobson argument~\cite{Jacobson2016} is UV-agnostic. We
provide a specific, non-holographic UV completion: a lattice of
self-modeling quantum systems whose interaction Hamiltonian is forced by
the algebraic structure of quantum theory itself.

The main result can be stated as follows. Self-modeling is the sole
physical premise. The lattice topology $G = (V, E)$ and the sign of
the coupling ($J > 0$, antiferromagnetic) are inputs. The Wilsonian
continuum limit is standard methodology. We present two routes to
Einstein's equations: an exact derivation via Jacobson's entanglement
equilibrium in the conformal regime, and a uniqueness argument via
Lovelock's theorem that extends to all $d \geq 2$ under the assumption
that entanglement perturbations yield a tensorial geometric equation.

The maximal vacuum entanglement hypothesis---that the vacuum maximizes
entanglement entropy at fixed stress-energy expectation
value---traditionally appears as an additional assumption in Jacobson's
framework. We argue that in the present context it is a \emph{structural
identification}: in the absence of pre-existing geometry, the modular
flow of the self-modeling lattice state defines the emergent geometric
flow, following the thermal time hypothesis of Connes and
Rovelli~\cite{ConnesRovelli1994}. The state whose modular flow defines a
smooth emergent geometry necessarily satisfies MVEH---the ``vacuum'' is
the geometry-defining state.  This is the only available identification
in a pre-geometric framework, though we do not prove existence or
uniqueness of such a state. We note that Sorce~\cite{Sorce2024} has shown
that geometric modular flow requires conformal symmetry; this caveat is
addressed in Sec.~\ref{sec:equilibrium}, where we observe that the
$SU(n)$ covariance of the self-modeling lattice places its infrared
description in the class of conformal field theories for which the
identification is clean.

This paper does not derive Newton's constant $\Gnewton$ (which depends on
the lattice entropy density~$\eta$ via $\Gnewton = 1/(4\eta)$), the
spacetime dimension (which is set by the lattice topology), or the
cosmological constant~$\Lambda$ (which appears as an undetermined
integration constant).  The sign chain producing attractive gravity
assumes the null energy condition (NEC).

The complete derivation chain, from self-modeling to Einstein's equations,
is summarized in Table~\ref{tab:chain}.

\begin{table*}
\caption{Derivation chain from self-modeling to Einstein's equations.
``Status'' gives the logical standing of each link. ``Novel?'' indicates
what is new in this work versus prior results. L7 is reframed as
a structural identification via the Connes--Rovelli thermal time
hypothesis~\cite{ConnesRovelli1994}, not an independent assumption.
\label{tab:chain}}
\begin{ruledtabular}
\begin{tabular}{clllll}
Link & Statement & Status & Section & Novel? & Key input \\ \hline
L1 & Self-modeling forces $M_n(\mathbb{C})^{sa}$ with L\"uders product
   & Proven & --- & Ref.~\cite{Paper5} & None \\
L2 & Lattice Hamiltonian $H = \sum J\,\swapop$ forced by $U(n)$ covariance
   & Derived & \ref{sec:lattice} & This work & Lattice topology, $J>0$ (inputs) \\
L3 & Lieb-Robinson velocity $\vlr$ gives emergent causal structure
   & Derived & \ref{sec:lattice} & This work & L2 \\
L4 & Area-law entanglement entropy
   & Established & \ref{sec:arealaw} & This work & L2, L3 \\
L5 & Entanglement first law: $\delta\ent = \delta\langle\modH\rangle$
   & Exact identity & \ref{sec:arealaw} & Standard QI & None \\
L6 & Wilsonian continuum limit
   & Standard methodology & \ref{sec:einstein} & Standard & L2 \\
L7 & Modular flow defines geometry (thermal time hypothesis)
   & Structural ident.\ & \ref{sec:equilibrium} & Reframed
   & Connes--Rovelli~\cite{ConnesRovelli1994} \\
L8a & $\Gab + \Lambda\,\gab = 8\pi\Gnewton\,\Tab$ (Jacobson route)
   & Derived (conformal) & \ref{sec:route-a} & Jacobson~\cite{Jacobson2016} applied
   & L1--L7, conformal $\modH$ \\
L8b & $\Gab + \Lambda\,\gab = 8\pi\Gnewton\,\Tab$ (Lovelock route)
   & Uniqueness ($d \geq 2$) & \ref{sec:route-b} & This work
   & L1--L7, Lovelock~\cite{Lovelock1971} \\
\end{tabular}
\end{ruledtabular}
\end{table*}

The novel contribution of this work is the assembly of Links~L1--L5 as
a UV completion, with L2--L4 being new results: the forced Hamiltonian
(L2), emergent causal structure (L3), and area-law bounds (L4).  L1 is
established in Paper~5~\cite{Paper5} and L5 is a standard identity of
quantum information theory.  We present two routes from L1--L7 to
Einstein's equations: Route~A applies Jacobson's entanglement
equilibrium argument~\cite{Jacobson2016} (exact in the conformal
regime), and Route~B uses Lovelock's uniqueness
theorem~\cite{Lovelock1971} (valid in all $d \geq 2$ spatial
dimensions without requiring conformal symmetry).  L7 is reframed from
an assumption to a structural identification.

The paper is organized as follows. Section~\ref{sec:lattice} defines the
self-modeling lattice: the local algebra $M_n(\mathbb{C})^{sa}$ at each
site, the SWAP interaction Hamiltonian forced by diagonal $U(n)$
covariance, and the emergent causal structure from the Lieb-Robinson
velocity~\cite{LiebRobinson1972, NachtergaeleSims2006}.
Section~\ref{sec:arealaw} establishes sub-volume entanglement scaling:
the WVCH thermal mutual information bound~\cite{WVCH2008}, known
ground-state entanglement properties of the Heisenberg
model~\cite{CalabreseCardy2004, Hastings2007}, and the entanglement
first law combined with modular Hamiltonian locality. Section~\ref{sec:equilibrium} presents the
thermal time hypothesis and its application to the self-modeling lattice,
recasting the MVEH assumption as a structural identification.
Section~\ref{sec:einstein} derives Einstein's equations via two routes:
Jacobson's entanglement equilibrium argument~\cite{Jacobson2016} in the
conformal regime (Route~A), and Lovelock's uniqueness
theorem~\cite{Lovelock1971} in $d \geq 2$ (Route~B).
Section~\ref{sec:numerical} presents numerical verification on small
lattices ($N = 8$--$20$): exact diagonalization benchmarks, area-law
scaling, modular Hamiltonian locality, and a qualitative MVEH check.
Section~\ref{sec:discussion} identifies remaining gaps, compares with
related work, and outlines future directions.
